\documentclass[11pt]{article}

\input{../preamble}


\begin{document}
	
\begin{center}
\Large{\textbf{CS 330, Fall 2026, Homework 4 \\
Due Wednesday, September 30, 2026 \\
}}
\end{center}

\emph{For full credit, you must follow the guidelines laid out on the first page of homework 1.}\\


\begin{problem}
Elvish languages of Middle-Earth (10 points)
\end{problem}

J.R.R. Tolkien, the author of the famous trilogy Lord of the Rings, was also a professor of philology (the study of language in oral and historical sources). Along with his novels he developed the mythology and origin of Middle-Earth and constructed multiple of its languages~\footnote{\url{https://en.wikipedia.org/wiki/Languages_constructed_by_Tolkien}}. He created the entire Elvish language family with 15 languages, grammar and vocabulary. The most famous of this is the inscription in the One Ring

\begin{center}
 \includegraphics[width=.39\textwidth]{One_Ring_inscription.png}
\end{center}

Recently Tolkien scholars have discovered a so far unknown manuscript which they suspect contains the description of another new language. The language consist of characters never seen by the scholars. There is a list of strings, a.k.a. ``words'' that the scholars believe are written in alphabetical order. Help the scholars decipher what the alphabet might look like and what the alphabetical order among the characters might be.\\

We call the discovered language \emph{developed} if we can indeed infer an alphabetical order from the manuscript (note that there may be multiple valid orders). If we cannot, then the language is \emph{incomplete}.

\textbf{Example 1:} a $\mathit{developed}$ language.
\begin{itemize}
    \item $\$\%\%*$
    \item $\&\&\#\&$
    \item $\&\&\#@$
\end{itemize}
One possible ordering of characters is $\$,\&,@,\%, *, \#$. (Check  for yourself that the words are indeed in alphabetical order.)\\

\textbf{Example 2:} an  $\mathit{incomplete}$ language.
\begin{itemize}
    \item $\$\%\%*$
    \item $\&\&\#\&$
    \item $\&\&\#\$$
\end{itemize}
The second example depicts an invalid list. This is because, if the first two words are in sorted order, then we must have $\$$ before $\&$; but if the final two words are in sorted order, then we must have $\&$ before $\$$, which is a contradiction. 

Note that you are not guaranteed that every letter will appear in the first position (refer to the examples above).\\

\begin{enumerate}
    \item This problem can be represented as a graph problem.  Write an algorithm that takes as input the manuscript $M$.  $M$ is a  nested \emph{array}, such that $M[i]$ consists of the character array of word $i$, e.g.  the first entry in example 1 is $M[0] = [\$,\%,\%,*]$. You may assume that all words are unique.  The output should be the graph adjacency list $G$ as a nested hash table.  $M$ consists of $W$ words and the language consists of $L$ distinct characters. For simplicity we assume that the length of each word is the same constant $c$. (Note, that it is possible that this graph is disconnected, or even may contain singleton nodes.)\\
    \item Write an algorithm that takes as input $G$ and either returns  the characters in alphabetical order if such an order exists, or return ``no alphabetical order''. If there are multiple possible orders, than any one of those is acceptable.

\end{enumerate}




\begin{problem}
Trip planning
\end{problem}

You are planning a vacation in a faraway country, and are looking at the train schedules in that country. 
The train schedule $train$ is made up of a set of $n$ $4$-tuples.
\begin{align*}
    train(i)=(depcity(i),arrcity(i),dep(i),arr(i))
\end{align*}
In the above schedule, $depcity$ and $arrcity$ are the cities that train $i$ leaves from and arrives in, 
respectively, and $dep$ and $arr$ represent the time that it does those two things. Basically, train $i$ 
leaves $depcity(i)$ at time $dep(i)$, and arrives in $arrcity(i)$ at time $arr(i)$.




Suppose you want to start at time $0$ in a particular city $s$, and wanted to see how quickly you could get to 
each city from your starting city. Write an algorithm that takes as input the train schedule $train$ and the starting city $s$ and efficiently finds the earliest arrival times to each city. Explain why it is correct.

Remember that this is a train schedule, and so you cannot catch a train if you arrive later at a city 
than it leaves.



\end{document}